\chapter{CƠ SỞ LÝ THUYẾT VÀ CÔNG TRÌNH LIÊN QUAN}
\label{ch:co_so}

\section{Bài toán phát hiện đối tượng}

\subsection{Phát biểu bài toán}

Cho một ảnh $I \in \mathbb{R}^{H \times W \times 3}$, bài toán phát hiện đối tượng yêu cầu sinh ra
một tập hợp
\begin{equation}
\label{eq:phatbieu}
\mathcal{D} = \{(c_i, b_i, s_i)\}_{i=1}^{N},
\end{equation}
trong đó $c_i \in \{1,\dots,K\}$ là nhãn lớp, $b_i = (x_i, y_i, w_i, h_i)$ là hộp bao và
$s_i \in [0,1]$ là điểm tin cậy. Số lượng vật thể $N$ không biết trước và thay đổi theo từng ảnh
--- đây chính là điểm khiến bài toán khó hơn phân lớp ảnh, nơi đầu ra luôn là một véc-tơ có kích
thước cố định.

Sự khác biệt về bản chất này dẫn tới hai chiến lược thiết kế. Chiến lược thứ nhất là biến bài toán
số lượng thay đổi thành bài toán số lượng cố định bằng cách rải dày các vị trí ứng viên trên ảnh
--- đó là con đường của các bộ phát hiện dựa trên hộp neo và lưới. Chiến lược thứ hai là giữ
nguyên bản chất tập hợp và học một phép gán --- đó là con đường của DETR.

\subsection{Định dạng chú thích}

Ba định dạng chú thích phổ biến được dùng trong đề tài này. Định dạng YOLO lưu mỗi vật thể trên
một dòng văn bản gồm \texttt{class\_id}, tâm hộp $(x_c, y_c)$ và kích thước $(w, h)$, tất cả đã
chuẩn hoá về đoạn $[0,1]$ theo kích thước ảnh. Định dạng COCO~\cite{lin2014coco} dùng một tệp JSON
duy nhất chứa ba khối \texttt{images}, \texttt{annotations} và \texttt{categories}, với hộp bao
lưu theo góc trên trái và kích thước tuyệt đối tính bằng điểm ảnh. Định dạng PASCAL
VOC~\cite{everingham2010pascalvoc} dùng tệp XML riêng cho mỗi ảnh.

Phép chuyển từ YOLO sang COCO là một bước bắt buộc trong đề tài vì nhánh DETR làm việc theo lược
đồ COCO. Công thức chuyển đổi:
\begin{align}
\label{eq:yolo2coco}
x_{\text{trái}} &= \left(x_c - \tfrac{w}{2}\right) \cdot W_{\text{ảnh}}, &
y_{\text{trên}} &= \left(y_c - \tfrac{h}{2}\right) \cdot H_{\text{ảnh}}, \\
w_{\text{tuyệt đối}} &= w \cdot W_{\text{ảnh}}, &
h_{\text{tuyệt đối}} &= h \cdot H_{\text{ảnh}}. \nonumber
\end{align}
Trường \texttt{area} bằng tích $w_{\text{tuyệt đối}} \cdot h_{\text{tuyệt đối}}$ và được bộ đánh
giá COCO dùng để phân nhóm vật thể thành nhỏ, trung bình và lớn khi tính mAP theo kích thước.

\section{Họ kiến trúc một giai đoạn}

\subsection{Nguyên lý chung}

Bộ phát hiện một giai đoạn bỏ hẳn bước đề xuất vùng. Ảnh đi qua một mạng xương sống tích chập để
sinh bản đồ đặc trưng, sau đó một đầu dự đoán gắn trực tiếp lên bản đồ đó sẽ xuất ra hộp bao và
điểm phân lớp cho từng vị trí. Vì toàn bộ quá trình chỉ cần một lượt truyền xuôi, tốc độ suy luận
cao hơn đáng kể so với dòng hai giai đoạn~\cite{redmon2016yolo}.

Điểm yếu cố hữu của cách làm này là mất cân bằng cực đoan giữa số vị trí nền và số vị trí chứa vật
thể. Focal Loss~\cite{lin2017focal} giải quyết bằng cách giảm trọng số đóng góp của các mẫu dễ:
\begin{equation}
\label{eq:focal}
\mathcal{L}_{\text{focal}}(p_t) = -\alpha_t \left(1 - p_t\right)^{\gamma} \log(p_t),
\end{equation}
trong đó $p_t$ là xác suất mô hình gán cho lớp đúng và $\gamma$ điều khiển mức độ giảm trọng số.

\subsection{Đa tỉ lệ và vật thể nhỏ}

Một biển báo cách xe 50 mét và một biển báo cách 5 mét có kích thước trên ảnh chênh nhau hàng chục
lần. Mạng kim tự tháp đặc trưng (FPN)~\cite{lin2017fpn} xử lý điều này bằng cách kết hợp đặc trưng
ở nhiều tầng phân giải: tầng nông giữ chi tiết không gian tốt cho vật thể nhỏ, tầng sâu giàu ngữ
nghĩa hơn cho vật thể lớn.

Dù vậy, vật thể nhỏ vẫn là điểm yếu dai dẳng~\cite{tong2020smallobject}. Nguyên nhân nằm ở bước
sải: đầu dự đoán P3 của YOLOv8 hoạt động ở bước sải 8, nghĩa là một vật thể 16 điểm ảnh chỉ trải
trên hai ô lưới. Khi ảnh đầu vào bị thu nhỏ để khớp kích thước huấn luyện, tình hình còn xấu hơn.
Đây là lý do vì sao chỉ số $AP$ tách riêng cho nhóm vật thể nhỏ có giá trị chẩn đoán cao hơn hẳn
mAP tổng, như sẽ thấy ở Chương~\ref{ch:ket_qua}.

\subsection{YOLOv8 và các hàm mất mát}

YOLOv8~\cite{jocher2023ultralytics} là một bộ phát hiện không dùng hộp neo. Hàm mất mát tổng gồm
ba thành phần:
\begin{equation}
\label{eq:yololoss}
\mathcal{L} = \lambda_{\text{box}} \mathcal{L}_{\text{CIoU}}
            + \lambda_{\text{cls}} \mathcal{L}_{\text{BCE}}
            + \lambda_{\text{dfl}} \mathcal{L}_{\text{DFL}}.
\end{equation}

Thành phần thứ nhất dùng CIoU~\cite{zheng2020diou}, một mở rộng của IoU có tính thêm khoảng cách
giữa hai tâm hộp và độ lệch tỉ lệ khung, nhờ đó vẫn cung cấp gradient hữu ích ngay cả khi hai hộp
chưa giao nhau --- trường hợp mà IoU thuần tuý cho gradient bằng không. Thành phần thứ ba là
Distribution Focal Loss~\cite{li2020gfl}, coi mỗi cạnh hộp bao là một phân phối rời rạc thay vì
một giá trị điểm, giúp việc hồi quy biên ổn định hơn khi ranh giới vật thể mờ.

Sau khi suy luận, bước hậu xử lý NMS loại bỏ các hộp trùng cho cùng một vật thể: các hộp được sắp
theo điểm tin cậy giảm dần, hộp có điểm cao nhất được giữ lại và mọi hộp còn lại có IoU với nó
vượt ngưỡng sẽ bị loại.

\section{Họ kiến trúc transformer}

\subsection{Cơ chế chú ý}

Transformer~\cite{vaswani2017attention} thay thế phép tích chập cục bộ bằng cơ chế chú ý cho phép
mọi vị trí tương tác trực tiếp với mọi vị trí khác:
\begin{equation}
\label{eq:attention}
\text{Attention}(Q, K, V) = \text{softmax}\!\left(\frac{QK^{\top}}{\sqrt{d_k}}\right) V .
\end{equation}
Hệ số $\sqrt{d_k}$ giữ cho phương sai của tích vô hướng không tăng theo số chiều, tránh đẩy hàm
softmax vào vùng bão hoà. Vision Transformer~\cite{dosovitskiy2021vit} chứng minh kiến trúc này áp
dụng được cho ảnh, nhưng cũng cho thấy nó cần lượng dữ liệu lớn hơn hẳn CNN --- một quan sát rất
quan trọng đối với đề tài này, vì nó dự báo trước điều gì sẽ xảy ra khi huấn luyện DETR trên bộ dữ
liệu chỉ 6.012 hộp bao.

Lý do của hiện tượng đó là \emph{thiên kiến quy nạp}. Mạng tích chập được cài sẵn giả định về tính
địa phương và bất biến tịnh tiến; transformer không có những giả định này và phải học chúng từ dữ
liệu. Khi dữ liệu dồi dào, việc không bị ràng buộc là lợi thế; khi dữ liệu khan hiếm, đó là bất
lợi.

\subsection{DETR: dự đoán theo tập hợp}

DETR~\cite{carion2020detr} đưa ra ba thay đổi so với các bộ phát hiện truyền thống. Thứ nhất, mô
hình dùng một tập cố định gồm 100 truy vấn đối tượng học được, mỗi truy vấn đề xuất tối đa một vật
thể. Thứ hai, việc ghép cặp giữa dự đoán và nhãn thật được giải bằng thuật toán
Hungarian~\cite{kuhn1955hungarian} với chi phí ghép cặp
\begin{equation}
\label{eq:hungarian}
\mathcal{L}_{\text{ghép}} = -\mathbb{1}_{\{c_i \neq \varnothing\}} \hat{p}_{\sigma(i)}(c_i)
  + \mathbb{1}_{\{c_i \neq \varnothing\}} \mathcal{L}_{\text{box}}(b_i, \hat{b}_{\sigma(i)}),
\end{equation}
trong đó $\mathcal{L}_{\text{box}}$ kết hợp khoảng cách $L_1$ và GIoU. Thứ ba, vì phép ghép là
song ánh nên mỗi nhãn thật khớp đúng một dự đoán, do đó mô hình không sinh hộp trùng và
\emph{không cần NMS}.

Cái giá của thiết kế này là tốc độ hội tụ. Bản gốc cần khoảng 300 đến 500 chu kỳ trên COCO. Nguyên
nhân được cho là phép gán Hungarian thay đổi mạnh giữa các bước đầu huấn luyện, khiến tín hiệu
giám sát không ổn định. Deformable DETR~\cite{zhu2021deformabledetr} khắc phục bằng cơ chế chú ý
chỉ lấy mẫu một số ít điểm tham chiếu, giảm số chu kỳ cần thiết xuống khoảng một phần mười. Gần
đây hơn, RT-DETR~\cite{zhao2024rtdetr} cho thấy dòng transformer có thể đạt tốc độ thời gian thực
và vượt YOLO ở cùng ngân sách --- nhưng vẫn với điều kiện được huấn luyện đầy đủ.

\section{Hệ thống độ đo}

\subsection{IoU, Precision và Recall}

Độ đo nền tảng là tỉ số giao trên hợp giữa hộp dự đoán $B_p$ và hộp thật $B_g$:
\begin{equation}
\label{eq:iou}
\text{IoU}(B_p, B_g) = \frac{|B_p \cap B_g|}{|B_p \cup B_g|} .
\end{equation}
Một dự đoán được tính là dương tính đúng khi lớp khớp và IoU vượt ngưỡng cho trước. Từ đó suy ra
\begin{equation}
\label{eq:pr}
\text{Precision} = \frac{TP}{TP + FP}, \qquad
\text{Recall} = \frac{TP}{TP + FN}.
\end{equation}

\subsection{mAP}

Sắp toàn bộ dự đoán theo điểm tin cậy giảm dần và tính cặp $(\text{Precision}, \text{Recall})$ tại
từng ngưỡng cho ta đường cong PR. Độ chính xác trung bình $AP$ là diện tích dưới đường cong này,
và $mAP$ là trung bình của $AP$ trên toàn bộ các lớp:
\begin{equation}
\label{eq:map}
AP = \int_{0}^{1} p(r)\, dr, \qquad
mAP = \frac{1}{K} \sum_{k=1}^{K} AP_k .
\end{equation}
Quy ước COCO báo cáo hai biến thể: mAP@0,5 tính tại một ngưỡng IoU duy nhất là 0,5, và mAP@0,5:0,95
lấy trung bình trên mười ngưỡng từ 0,5 đến 0,95 với bước 0,05. Biến thể thứ hai khắt khe hơn nhiều
vì nó đòi hỏi hộp bao phải khít, chứ không chỉ chồng lấn ở mức tối thiểu.

Một hệ quả quan trọng của việc mAP là độ đo \emph{xếp hạng}: khi tính mAP, không được lọc dự đoán
theo ngưỡng tin cậy trước, vì làm vậy sẽ cắt mất phần đuôi của đường cong PR và bóp méo kết quả.
Ngưỡng tin cậy chỉ nên áp dụng ở khâu hiển thị.

\subsection{Độ đo chi phí triển khai}

Ba đại lượng bổ sung được dùng để phản ánh ràng buộc thực tế: độ trễ suy luận tính bằng
mili-giây cho mỗi ảnh, tốc độ khung hình bằng nghịch đảo của độ trễ, và dung lượng mô hình tính
bằng megabyte. Cả ba đều phụ thuộc phần cứng nên phải luôn báo cáo kèm thiết bị đo. Ngoài ra, do
GPU thực thi bất đồng bộ, phép đo thời gian phải có bước đồng bộ hoá tường minh, nếu không sẽ ghi
nhận thiếu thời gian tính toán thực.

\section{Nén mô hình}

\subsection{Lượng tử hoá}

Lượng tử hoá thay biểu diễn dấu phẩy động 32 bit bằng biểu diễn có độ rộng bit nhỏ hơn. Phép ánh
xạ affine sang số nguyên 8 bit~\cite{jacob2018quantization} có dạng
\begin{equation}
\label{eq:quant}
r \approx S\,(q - Z),
\end{equation}
trong đó $r$ là giá trị thực, $q$ là giá trị nguyên, $S$ là hệ số tỉ lệ và $Z$ là điểm không. Ở
mức FP16, việc giảm độ rộng bit gần như không ảnh hưởng tới chất lượng mô hình trong nhiều tác
vụ~\cite{micikevicius2018mixedprecision}, trong khi INT8 nén mạnh hơn nhưng có thể mất độ chính
xác và --- điều thường bị bỏ qua --- chỉ thực sự nhanh hơn nếu môi trường thực thi có nhân tính
toán số nguyên được tối ưu.

\subsection{Chưng cất tri thức}

Chưng cất tri thức~\cite{hinton2015distillation} huấn luyện một mô hình học trò nhỏ bắt chước
phân phối đầu ra mềm của một mô hình thầy lớn hơn. Hàm mất mát kết hợp
\begin{equation}
\label{eq:distill}
\mathcal{L} = (1-\alpha)\,\mathcal{L}_{\text{cứng}}
            + \alpha\, T^{2}\, \mathcal{L}_{\text{KL}}\!\left(\sigma(z_s/T) \,\|\, \sigma(z_t/T)\right),
\end{equation}
với $T$ là nhiệt độ làm mềm phân phối và $z_s, z_t$ là logit của học trò và thầy. Giả thuyết nền
là phân phối mềm của thầy mang thông tin về quan hệ giữa các lớp mà nhãn cứng không có. Tuy nhiên
lợi ích này không được bảo đảm: khi khoảng cách năng lực giữa thầy và trò quá nhỏ, hoặc khi bản
thân thầy chưa đủ tốt, tín hiệu chưng cất có thể trở thành nhiễu.

\section{Công trình liên quan về biển báo giao thông}

GTSRB~\cite{stallkamp2012gtsrb} là bộ chuẩn kinh điển cho \emph{nhận dạng} biển báo, gồm các ảnh
đã cắt sát quanh từng biển; đây là bài toán phân lớp chứ không phải phát hiện.
GTSDB~\cite{houben2013gtsdb} bổ sung phần \emph{phát hiện} với ảnh đường phố đầy đủ.
TT100K~\cite{zhu2016tt100k} mở rộng quy mô sang bối cảnh Trung Quốc với ảnh toàn cảnh độ phân giải
cao, còn Mapillary Traffic Sign~\cite{ertler2020mapillary} cung cấp dữ liệu đa quốc gia thu từ
camera hành trình.

Phân biệt giữa bộ dữ liệu kiểu nhận dạng và bộ dữ liệu kiểu phát hiện có ý nghĩa trực tiếp với đề
tài này: như sẽ trình bày ở Mục~\ref{sec:thanhphan}, bộ dữ liệu được dùng là một tập trộn chứa cả
hai kiểu, và chính điều đó chi phối cách diễn giải kết quả. Hiện tượng bộ dữ liệu mang thiên lệch
đặc thù và mô hình học phải thiên lệch đó thay vì học đặc trưng tổng quát đã được chỉ ra từ
lâu~\cite{torralba2011bias}, còn việc chỉ số trên tập kiểm tra không chuyển được sang dữ liệu mới
cũng đã được ghi nhận trên các bộ chuẩn lớn~\cite{recht2019imagenetv2}.
